\documentclass{beamer}

\usetheme{Madrid}
\usecolortheme{seahorse}
\usepackage[utf8]{inputenc}
\usepackage[spanish]{babel}
\usepackage{listings}
\usepackage{xcolor}

% --- Configuración del estilo del código C# ---
\lstdefinestyle{CSharpStyle}{
  language=[Sharp]C,
  basicstyle=\scriptsize\ttfamily,
  keywordstyle=\color{blue}\bfseries,
  stringstyle=\color{teal},
  commentstyle=\color{gray}\itshape,
  showstringspaces=false,
  breaklines=true,
  frame=single,
  rulecolor=\color{gray!30},
  tabsize=2,
  captionpos=b,
  numbers=left,
  numberstyle=\tiny\color{gray},
}

\title{Uso de Stack (Pila) en C\# para el Desarrollo de Videojuegos con Unity}
\author{COM200-Programacion Avanzada}
\date{\today}

\begin{document}

\begin{frame}
\titlepage
\end{frame}

%------------------------------------------------
\begin{frame}{Concepto general}
\begin{itemize}
    \item Un \textbf{Stack (pila)} en C\# es una estructura de datos del tipo \textbf{LIFO (Last In, First Out)}.
    \item El último elemento que entra es el primero en salir.
    \item En un videojuego, puede imaginarse como una \textbf{torre de platos apilados}.
    \item En Unity, esta estructura es útil para:
    \begin{itemize}
        \item Recolectar y entregar objetos apilados (platos, cajas, recursos).
        \item Registrar acciones recientes o implementar “undo”.
        \item Controlar inventarios tipo pila.
    \end{itemize}
\end{itemize}
\end{frame}

%------------------------------------------------
\begin{frame}[fragile]{Declaración de un Stack}
Antes de usar una pila en C\#, se importa el espacio de nombres:
\begin{lstlisting}
using System.Collections.Generic;
\end{lstlisting}

Luego, se declara la pila del tipo deseado:
\begin{lstlisting}
Stack<GameObject> pilaDePlatos = new Stack<GameObject>();
\end{lstlisting}
\end{frame}

%------------------------------------------------
\begin{frame}[fragile]{Operación Push() — Agregar elementos}
\textbf{Push()} agrega un elemento en la parte superior de la pila.
\begin{lstlisting}
pilaDePlatos.Push(plato);
\end{lstlisting}

Ejemplo práctico en Unity:
\begin{lstlisting}
public class RecolectorDePlatos : MonoBehaviour
{
    private Stack<GameObject> pilaDePlatos = new Stack<GameObject>();

    public void RecogerPlato(GameObject plato)
    {
        pilaDePlatos.Push(plato);
        Debug.Log($"Plato {plato.name} añadido a la pila. Total: {pilaDePlatos.Count}");
    }
}
\end{lstlisting}
\end{frame}

%------------------------------------------------
\begin{frame}[fragile]{Operación Pop() — Quitar y obtener el último elemento}
\textbf{Pop()} elimina el último elemento agregado y lo devuelve.
\begin{lstlisting}
GameObject ultimoPlato = pilaDePlatos.Pop();
Debug.Log($"Plato entregado: {ultimoPlato.name}");
\end{lstlisting}

Validación antes de eliminar:
\begin{lstlisting}
if (pilaDePlatos.Count > 0)
{
    GameObject plato = pilaDePlatos.Pop();
}
\end{lstlisting}
\end{frame}

%------------------------------------------------
\begin{frame}[fragile]{Operación Peek() — Consultar sin eliminar}
Permite ver el último elemento de la pila sin eliminarlo.
\begin{lstlisting}
GameObject platoSuperior = pilaDePlatos.Peek();
Debug.Log("El plato superior es: " + platoSuperior.name);
\end{lstlisting}

Ideal para mostrar en la interfaz qué plato será entregado a continuación.
\end{frame}

%------------------------------------------------
\begin{frame}[fragile]{Otras operaciones importantes}
\begin{itemize}
    \item \textbf{Clear()}: Vacía toda la pila.
    \begin{lstlisting}
    pilaDePlatos.Clear();
    \end{lstlisting}

    \item \textbf{Count}: Devuelve la cantidad de elementos.
    \begin{lstlisting}
    int cantidad = pilaDePlatos.Count;
    \end{lstlisting}

    \item \textbf{Contains()}: Comprueba si un elemento está presente.
    \begin{lstlisting}
    if (pilaDePlatos.Contains(platoEspecial))
        Debug.Log("¡El plato dorado está en la pila!");
    \end{lstlisting}
\end{itemize}
\end{frame}

%------------------------------------------------
\begin{frame}[fragile]{Ejemplo completo: Recolector de platos}
\begin{lstlisting}
using UnityEngine;
using System.Collections.Generic;

public class RecolectorDePlatos : MonoBehaviour
{
    private Stack<GameObject> pilaDePlatos = new Stack<GameObject>();
    public Transform puntoDeApilado;
    private float alturaEntrePlatos = 0.2f;

    public void RecogerPlato(GameObject plato)
    {
        Vector3 nuevaPos = puntoDeApilado.position + 
                           Vector3.up * (pilaDePlatos.Count * alturaEntrePlatos);

        GameObject nuevoPlato = Instantiate(plato, nuevaPos, 
                                            Quaternion.identity, puntoDeApilado);
        pilaDePlatos.Push(nuevoPlato);
    }

    public void EntregarPlato()
    {
        if (pilaDePlatos.Count > 0)
        {
            GameObject platoSuperior = pilaDePlatos.Pop();
            Destroy(platoSuperior);
        }
    }
}
\end{lstlisting}
\end{frame}

%------------------------------------------------
\begin{frame}{Ventajas de usar Stack en Unity}
\begin{itemize}
    \item Mantiene un orden lógico inverso (LIFO).
    \item Facilita la gestión de objetos apilados o reversibles.
    \item Simplifica la programación de mecánicas de recolección y entrega.
    \item Mejora la claridad y la eficiencia del código.
\end{itemize}
\end{frame}

%------------------------------------------------
\begin{frame}{Comparación con otras estructuras de datos}
\begin{table}[]
\centering
\begin{tabular}{|l|l|l|}
\hline
\textbf{Estructura} & \textbf{Orden de acceso} & \textbf{Uso común en videojuegos} \\ \hline
Stack & LIFO (último entra, primero sale) & Recolector de platos, historial \\ \hline
Queue & FIFO (primero entra, primero sale) & Turnos, spawns secuenciales \\ \hline
List & Acceso directo por índice & Inventarios, listas de objetos \\ \hline
\end{tabular}
\end{table}
\end{frame}

%------------------------------------------------
\begin{frame}{Conclusión}
\begin{itemize}
    \item La clase \textbf{Stack} en C\# permite manejar datos en orden inverso.
    \item Es ideal para mecánicas donde el último elemento agregado debe procesarse primero.
    \item En Unity, facilita la implementación de sistemas de recolección, almacenamiento y entrega.
    \item Su simplicidad y claridad la convierten en una herramienta valiosa para el desarrollo de videojuegos.
\end{itemize}
\end{frame}

\end{document}
